\begin{abstract}
Este trabalho apresenta um estudo teórico e experimental sobre o Problema do Caminho Hamiltoniano em grafos não direcionados e conexos. Teoricamente, são apresentadas as definições formais do problema, demonstra-se seu pertencimento à classe NP e prova-se sua NP-completude por meio de reduções polinomiais a partir do 3-SAT, além de uma discussão sobre métodos para resolvê-lo. Experimentalmente, são comparadas duas abordagens exatas: uma baseada em programação dinâmica (Bellman--Held--Karp) e outra baseada em redução para SAT com um solucionador Conflict-Driven Clause Learning, avaliadas quanto ao tempo de execução e ao consumo de memória em instâncias de diferentes tamanhos e densidades. Os resultados indicam que a abordagem baseada em programação dinâmica apresenta melhor desempenho temporal para as instâncias consideradas, enquanto a redução para SAT sofre crescimento mais acentuado do tempo de execução; e que ambas apresentam crescimento exponencial no uso de memória, ainda que na redução para SAT seja mais suave.
\end{abstract}